\section{Búsqueda y Optimización}

\subsection{Búsqueda Binaria sobre la Respuesta + Two Pointers}

Preguntas guía — ¿cuándo usarlo?

\begin{itemize}
\item ¿La respuesta es un número y existe una condición monótona (si X funciona, todo valor mayor/menor también funciona)?
\item ¿Puedo verificar en O(n) o O(n log n) si un valor candidato "cumple" la condición pedida?
\item ¿Dentro de esa verificación necesito contar pares/elementos con una condición de suma sobre un arreglo ordenado? (ahí entra Two Pointers)
\end{itemize}

¿Para qué sirve? Cuando no es fácil calcular la respuesta directamente pero sí verificar si un candidato es válido, y esa validez es monótona, se "busca" la respuesta con búsqueda binaria sobre el rango de valores posibles. Two Pointers acelera el conteo en arreglos ordenados de O(n²) a O(n).

Complejidad: Búsqueda binaria: O(log(rango)) iteraciones. Cada verificación con Two Pointers: O(n). Total: O(n log(rango)).

Fundamento teórico: La búsqueda binaria sobre la respuesta funciona porque el conjunto de valores "válidos" forma un intervalo continuo (monotonía), reduciendo el espacio de búsqueda a la mitad en cada paso. Two Pointers explota que, en un arreglo ordenado, mover un extremo cambia la condición de suma de forma predecible, evitando recorrer todas las parejas.

Ejercicio aplicado

Dados N precios de categorías de joyas, un cliente que es el K-ésimo en la fila debe comprar dos joyas de categorías distintas sin repetir una combinación ya usada por otro cliente. Hallar el mínimo dinero que garantiza poder comprar: la K-ésima suma más pequeña entre todos los pares de categorías distintas.

E/S: N no está acotado explícitamente pero puede ser grande → lectura total + tokenización (patrón 2) para el arreglo de precios.

Código solución

\begin{lstlisting}[language=Python]
import sys
 
def count_pairs_leq(prices, n, limit):
    # cuenta cuantos pares cumplen arr[i] + arr[j] <= limit
    # usando dos punteros sobre un arreglo YA ORDENADO
    count = 0
    left, right = 0, n - 1
    while left < right:
        if prices[left] + prices[right] <= limit:
            count += right - left       # todos los pares con "left" fijo
            left += 1
        else:
            right -= 1
    return count
 
def solve():
    data = sys.stdin.buffer.read().split()
    n = int(data[0]); k = int(data[1])
    prices = list(map(int, data[2:2 + n]))
    prices.sort()
 
    low = prices[0] + prices[1]         # menor suma posible
    high = prices[n - 2] + prices[n - 1]  # mayor suma posible
 
    while low < high:
        mid = low + (high - low) // 2
        if count_pairs_leq(prices, n, mid) >= k:
            high = mid                    # mid cumple: intenta algo menor
\end{lstlisting}


\begin{lstlisting}[language=Python]
        else:
            low = mid + 1                # mid no cumple: sube el limite
 
    print(low)
 
solve()
\end{lstlisting}


\subsection{------NEED: Enumeración por Máscaras de Bits (Bitmask Brute Force)}

Preguntas guía — ¿cuándo usarlo?

\begin{itemize}
\item ¿Hay un pequeño número (≤ \textasciitilde{}20) de decisiones binarias desconocidas (bits, elecciones sí/no) y hace falta probar combinaciones?
\item ¿El espacio total de combinaciones (2\textasciicircum{}k) es lo bastante chico para probarlas todas dentro del límite de tiempo?
\item ¿Cada combinación se puede armar y evaluar rápido a partir de su representación binaria?
\end{itemize}

¿Para qué sirve? Representar y recorrer exhaustivamente todas las configuraciones posibles de un conjunto pequeño de decisiones binarias. Cada máscara de bits codifica una asignación distinta de valores, permitiendo generar candidatos de forma sistemática y verificar cuál satisface la condición del problema.

Complejidad: O(2\textasciicircum{}k · costo de evaluar una combinación), con k = número de incógnitas binarias.

Fundamento teórico: Cualquier asignación de k bits desconocidos corresponde biunívocamente a un entero entre 0 y 2\textasciicircum{}k−1. Recorriendo ese rango y extrayendo el bit j con (mascara \textgreater{}\textgreater{} j) \& 1, se reconstruye cada combinación posible sin necesidad de recursión ni backtracking explícito. Es la enumeración más directa de un espacio de búsqueda combinatorio pequeño.

Ejercicio aplicado

Un mensaje M' y un divisor N' se recibieron con hasta 16 bits en total marcados como perdidos ('*'). Reconstruir un mensaje M compatible (M completando los '*' de M') tal que exista un N (completando los '*' de N') donde N divide a M exactamente.

E/S: |M'| ≤ 500 pero el total de '*' (en M' y N' combinados) es ≤ 16 → lectura total con readline es más que suficiente; el costo real está en las 2\textasciicircum{}16 combinaciones, no en la entrada.

Código solución

\begin{lstlisting}[language=Python]
import sys

input = sys.stdin.readline


def generate(strings):
    pos = []            # posiciones donde existen bits desconocidos
    for i, c in enumerate(strings):
        if c == '*':
            pos.append(i)
    k = len(pos)

    ans = []
    for mask in range(1 << k):          # cada mascara representa una asignacion distinta de los k bits perdidos
        cur = strings[:]
        for j in range(k):
            if (mask >> j) & 1:
                cur[pos[j]] = '1'
            else:
                cur[pos[j]] = '0'
        ans.append(''.join(cur))
    return ans

M = list(input().strip())
N = list(input().strip())

messages = generate(M)
divisors = generate(N)

for m in messages:              # probar todas las combinaciones posibles
    value_m = int(m, 2)
    for n in divisors:
        value_n = int(n, 2)

        if value_n == 0:
            continue            # la division entre cero no existe
        if value_m % value_n == 0:        # la primera combinacion valida es respuesta
            print(m)
            sys.exit()
\end{lstlisting}